% SEGMENT OLP-0531-S01
% Part: set-theory
% Chapter: story
% Section: extensionality

\documentclass[../../../include/open-logic-section]{subfiles}

\begin{document}

\olfileid{sth}{story}{extensionality}
\olsection{உறுப்புசார் சமத்துவம்}

முதலில் சொல்ல வேண்டியது, கணங்கள் அவற்றின் !!{element}s ஆல்
தனிப்படுத்தப்படுகின்றன என்பதே. இன்னும் துல்லியமாக:

\begin{axiom}[உறுப்புசார் சமத்துவம்]
கணங்கள் $A$ மற்றும் $B$ ஒரே !!{element}s ஐக் கொண்டிருந்தால், $A$ மற்றும்
$B$ ஒரே கணம்.
\[
  \lforall[A][\lforall[B][(\lforall[x][(x \in A \liff x \in B)] \lif
  \eq[A][B])]]
\]
\end{axiom}

\olref[sfr][][]{part} முழுவதிலும் இதை நாம் எடுகோளாகக் கொண்டோம். ஆனால்
இதை மீண்டும் கூறுவது பயனுள்ளது. ஒரு கணம் அதன் !!{element}s ஆல்
நிர்ணயிக்கப்படுகிறது என்ற அடிப்படைக் கருத்தை உறுப்புசார் சமத்துவ அடிகோள்
வெளிப்படுத்துகிறது. (இதனால் கணங்களை \emph{கருத்துருக்களிலிருந்து}
வேறுபடுத்தலாம்: துல்லியமாக அதே பொருள்கள் பல வேறு கருத்துருக்களின் கீழ்
வரக்கூடும்.)

இந்தக் கொள்கையை ஏன் ஏற்க வேண்டும்? உறுப்புசார் சமத்துவத்தை மறுப்பது,
\emph{“கணக் கோட்பாடு”} எனக்கூட அழைக்கப்படக்கூடிய எதையும் கைவிடும் முடிவாகும்
என்று கூறுவது நம்பத்தக்கதே. கணக் கோட்பாடு என்பது உறுப்புகளால்
நிர்ணயிக்கப்படும் தொகுப்புகளின் கோட்பாடே; அதற்குக் குறையுமில்லை,
மிகையுமில்லை.

ஆனால் \olref[sth][][]{part} இல் உள்ள உண்மையான சவால், \emph{எந்தக் கணங்கள்
உள்ளன} என்பதைச் சொல்லும் கொள்கைகளை வகுப்பதாகும். இந்தக் கேள்விக்கு
உண்மையிலேயே “வெளிப்படையான” ஒரே விடை மெய்ப்பிக்கத்தக்கவாறு தவறானது
என்பது தெரியவருகிறது.

\end{document}
